\chapter{Balance, deuda y sostenibilidad}

\marginal{tres flujos, tres acervos, con su nombre}
Este capítulo empieza nombrando sus objetos, porque el error característico de la
literatura sobre deuda pública mexicana no es aritmético sino de nomenclatura: la
cifra correcta con el nombre equivocado.

Los tres flujos de 2025 son el \textbf{balance presupuestario}, un déficit de
1,170,566.5 millones de pesos o 3.2\% del PIB; el balance público; y los
\textbf{requerimientos financieros del sector público}, 1,428,348.1 millones o
3.9\% del PIB. Los tres acervos son el \textbf{saldo histórico de los
requerimientos financieros}, 18,591,005.5 millones o 51.4\% del PIB; la
\textbf{deuda neta del sector público federal}, 50.9\%; y la \textbf{deuda bruta
del sector público no financiero}, 54.9\%.

\section{Techos, endeudamiento y déficit no son lo mismo}

\begin{table}[htbp]
\centering
\caption{Techos, endeudamiento y déficit: tres objetos distintos}
\label{tab:C12_3}
\small
\begin{tabular}{lSSll}
\toprule
{Objeto} & {2024} & {2025} & {Fuente} & {Qué es} \\
\midrule
Endeudamiento neto interno del Gobierno Federal (techo) & \num{1990000.0} & \num{1580000.0} & ILIF art. 2o. & autorizacion, no gasto \\
Endeudamiento neto externo del sector publico (techo, millones de dolares) & \num{18000.0} & \num{15500.0} & ILIF art. 2o. & en dolares, no en pesos \\
Endeudamiento neto interno del Gobierno Federal (partida informativa) & \num{1906069.4} & \num{1576170.0} & LIF art. 1o. renglon 0.01.01 & lo que se preve ejercer \\
Deficit presupuestario & \num{1693000.0} & \num{1170566.5} & DEC art. 2 y CGPE & el flujo, distinto de los dos anteriores \\
\bottomrule
\end{tabular}
\par\vspace{2pt}\footnotesize Fuente: elaboración propia del ITED con la ILIF y la LIF aprobada, y los decretos. El renglón externo está en millones de dólares y no se suma con los demás.
\end{table}


Tres objetos que suelen confundirse y que en 2025 valen tres cosas distintas. El
\textbf{techo} de endeudamiento neto interno que autoriza el Congreso es de
1,580,000 millones de pesos: es una autorización, no un gasto. El
\textbf{endeudamiento informativo} que la Ley de Ingresos prevé ejercer es de
1,576,170.0. El \textbf{déficit presupuestario} es de 1,170,566.5. Ninguno es el
otro, y la diferencia entre el primero y el tercero, más de 400 mil millones de
pesos, es margen legal que no se piensa usar.

Los techos caen con fuerza: el interno 20.6\% nominal desde 1,990,000 y el externo
de 18,000 a 15,500 millones de dólares. \textbf{El renglón externo está en dólares
y no se suma con los de pesos}, advertencia que parece obvia y que el género
incumple con regularidad.

\section{El costo financiero}

El costo financiero es de \textbf{1,388,373.6 millones de pesos, 3.8\% del PIB}, y
crece 5.4\% real contra el aprobado de 2024 y 8.5\% contra el cierre estimado. Es
el ramo que más sube de todo el presupuesto: 82,550.6 millones de pesos de 2025.

\textbf{En un paquete que consolida dos puntos del PIB, el servicio de la deuda es
la partida que crece.} No hay contradicción: la consolidación reduce el déficit
del año, no el acervo acumulado, y la tasa real que se paga sobre ese acervo sigue
siendo alta. El capítulo siguiente lo cuantifica.

\section{La descomposición en cuatro variables}

\marginal{lo que el género no trae}
El género describe la deuda; no la descompone. Esta sección es la aportación
propia del documento y por eso pasa primero por una prueba de aceptación.

El marco es la identidad
\[
d_t = \frac{d_{t-1}}{1+n} + \mathrm{rfspf} + \mathrm{fx} + \text{otros},
\qquad n = (1+g)(1+\pi) - 1,
\]
con el efecto denominador repartido entre crecimiento e inflación por una
convención que se declara, y con la \textbf{compensación por inflación que el
sector público paga} contabilizada aparte, porque viaja dentro del costo
financiero y de las adecuaciones a registros presupuestarios. El efecto neto de la
inflación es la suma de los dos, no el primero solo.

\begin{table}[htbp]
\centering
\caption{Descomposición del cambio de la razón SHRFSPF/PIB en cuatro variables (puntos del PIB)}
\label{tab:C12_4}
\small
\begin{tabular}{lSSSSSSSSSSS}
\toprule
{Caso} & {$d_{t-1}$} & {$d_t$} & {Obs.} & {RFSPF} & {Crec.} & {Infl. denom.} & {Compens.} & {Infl. neta} & {Tipo cambio} & {Suma} & {Residuo} \\
\midrule
Prueba 2022, añada del CGPE 2024 (la validada en la rubrica 2023) & \num{49.2} & \num{47.7} & \num{-1.5} & \num{4.3} & \num{-1.73} & \num{-3.09} & \num{2.97} & \num{-0.12} & \num{-0.8} & \num{-1.32} & \num{-0.18} \\
Prueba 2022, añada que nombra la instruccion 2025 (50.7 -> 49.4) & \num{50.7} & \num{49.4} & \num{-1.3} & \num{4.5} & \num{-1.78} & \num{-3.18} & \num{3.06} & \num{-0.12} & \num{-0.82} & \num{-1.29} & \num{-0.01} \\
Aplicacion a 2025 (SHRFSPF 51.4 -> 51.4) & \num{51.4} & \num{51.4} & \num{0.0} & \num{3.9} & \num{-1.06} & \num{-2.12} & \num{2.07} & \num{-0.05} & \num{-0.79} & \num{-0.07} & \num{0.07} \\
\bottomrule
\end{tabular}
\par\vspace{2pt}\footnotesize Fuente: construcción propia del ITED con el CGPE 2025. Los dos primeros renglones son la prueba de aceptación del marco sobre 2022, con las dos añadas que estaban en disputa.
\end{table}


\subsection{Prueba de aceptación}

El marco se prueba sobre 2022, un año en que la razón cayó. \textbf{Hubo que
correrlo dos veces} porque las cifras de 2022 existen en más de una añada y no
coinciden: con la del CGPE 2024, la razón va de 49.2 a 47.7 con requerimientos de
4.3 puntos, y el marco reproduce una caída de 1.32 contra 1.5 observada, con un
residuo de 0.18. Con la añada alternativa, de 50.7 a 49.4 con 4.5 puntos, el marco
da 1.29 contra 1.3, con un residuo de 0.01.

\textbf{Las dos pasan, y la diferencia entre ellas es el mismo problema de añada
que este proyecto documenta en otros.} Se reportan las dos y no se elige una en
silencio.

\subsection{Aplicación a 2025}

El acervo no se mueve: 51.4\% del PIB en 2024 estimado y 51.4 en 2025. Que no se
mueva con un requerimiento de 3.9 puntos del PIB es el resultado que hay que
explicar, y el marco lo explica así:

\begin{itemize}
\item Los requerimientos financieros suman \textbf{$+3.9$} puntos.
\item El crecimiento real del PIB resta \textbf{$-1.06$}.
\item La inflación resta \textbf{$-2.12$} por el lado del denominador, pero el
sector público \textbf{paga $+2.07$} de compensación por inflación sobre la deuda
indexada. \textbf{El efecto neto de la inflación es de $-0.05$ puntos, es decir
prácticamente cero.}
\item La apreciación del peso proyectada, de 19.7 a 18.5 al cierre, resta
\textbf{$-0.79$} sobre los 12.9 puntos de deuda externa.
\item La suma da $-0.07$ y el cambio observado es 0.0: \textbf{un residuo de siete
centésimas de punto del PIB}.
\end{itemize}

\marginal{la inflación no diluye}
Dos conclusiones que no son de manual. La primera: \textbf{la inflación no
diluye la deuda cuando la tasa nominal la sigue}. Lo que aparenta ser una
licuación de dos puntos se cancela casi exactamente contra la compensación que se
paga, y ese pago está dentro del costo financiero que crece 5.4\% real. La segunda:
\textbf{el acervo se mantiene plano gracias a un supuesto de tipo de cambio}. Casi
ocho décimas de punto del PIB dependen de que el peso se aprecie a 18.5. Si
cerrara en 19.7, la razón subiría a 52.2 sin que nada de la política fiscal
cambiara.

\section{El escenario a 2030}

\footnotesize
\setlength{\tabcolsep}{4pt}
\begin{longtable}{L{3.7cm}ccccccccc}
\caption{Perspectivas de finanzas públicas, 2024--2030 (\textbackslash\{\}\% del PIB)}\label{tab:C12_2}\\
\toprule
{Concepto} & {2024 apr.} & {2024 est.} & {2025} & {2026} & {2027} & {2028} & {2029} & {2030} & {Dif. 25--30} \\
\midrule\endfirsthead
\toprule
{Concepto} & {2024 apr.} & {2024 est.} & {2025} & {2026} & {2027} & {2028} & {2029} & {2030} & {Dif. 25--30} \\
\midrule\endhead
\bottomrule\endfoot
{I. RFSP} & \num{-5.4} & \num{-5.9} & \num{-3.9} & \num{-3.2} & \num{-2.9} & \num{-2.9} & \num{-2.9} & \num{-2.9} & \num{1.1} \\
{II. Necesidades de financiamiento fuera del presupuesto} & \num{-0.5} & \num{-0.9} & \num{-0.7} & \num{-0.5} & \num{-0.5} & \num{-0.5} & \num{-0.5} & \num{-0.5} & \num{0.2} \\
{Pidiregas} & \num{-0.1} & \num{0.0} & \num{-0.1} & \num{-0.1} & \num{-0.1} & \num{-0.1} & \num{-0.1} & \num{-0.1} & \num{0.0} \\
{Requerimientos financieros del IPAB} & \num{-0.1} & \num{0.0} & \num{-0.1} & \num{-0.1} & \num{-0.1} & \num{-0.1} & \num{-0.1} & \num{-0.1} & \num{0.0} \\
{Requerimientos financieros del Fonadin} & \num{-0.1} & \num{0.0} & \num{0.0} & \num{0.0} & \num{0.0} & \num{0.0} & \num{0.0} & \num{0.0} & \num{0.1} \\
{Adecuaciones a registros presupuestarios} & \num{-0.2} & \num{-0.8} & \num{-0.4} & \num{-0.3} & \num{-0.3} & \num{-0.3} & \num{-0.3} & \num{-0.3} & \num{0.1} \\
{III. Balance presupuestario} & \num{-4.9} & \num{-5.0} & \num{-3.2} & \num{-2.7} & \num{-2.4} & \num{-2.4} & \num{-2.4} & \num{-2.4} & \num{0.9} \\
{III.A Ingresos presupuestarios} & \num{21.3} & \num{22.1} & \num{22.3} & \num{22.2} & \num{22.1} & \num{22.0} & \num{21.9} & \num{21.8} & \num{-0.5} \\
{Petroleros} & \num{3.0} & \num{3.1} & \num{3.2} & \num{3.1} & \num{3.1} & \num{3.0} & \num{3.0} & \num{3.0} & \num{-0.2} \\
{Petroleros, Gobierno Federal} & \num{0.8} & \num{0.6} & \num{0.8} & \num{0.9} & \num{0.9} & \num{0.9} & \num{0.9} & \num{0.9} & \num{0.1} \\
{Petroleros, Pemex} & \num{2.2} & \num{2.5} & \num{2.4} & \num{2.2} & \num{2.1} & \num{2.1} & \num{2.1} & \num{2.1} & \num{-0.3} \\
{No petroleros} & \num{18.3} & \num{19.0} & \num{19.1} & \num{19.1} & \num{19.0} & \num{18.9} & \num{18.9} & \num{18.8} & \num{-0.3} \\
{No petroleros, Gobierno Federal} & \num{15.1} & \num{15.5} & \num{15.7} & \num{15.7} & \num{15.6} & \num{15.5} & \num{15.5} & \num{15.4} & \num{-0.3} \\
{Tributarios} & \num{14.4} & \num{14.5} & \num{14.6} & \num{14.6} & \num{14.6} & \num{14.5} & \num{14.5} & \num{14.4} & \num{-0.2} \\
{No tributarios} & \num{0.8} & \num{1.0} & \num{1.0} & \num{1.0} & \num{1.0} & \num{1.0} & \num{1.0} & \num{1.0} & \num{0.0} \\
{Organismos y empresas} & \num{3.1} & \num{3.5} & \num{3.4} & \num{3.4} & \num{3.4} & \num{3.4} & \num{3.4} & \num{3.4} & \num{0.0} \\
{IMSS e ISSSTE} & \num{1.8} & \num{2.0} & \num{1.9} & \num{1.9} & \num{1.9} & \num{1.9} & \num{1.9} & \num{1.9} & \num{0.0} \\
{CFE} & \num{1.3} & \num{1.5} & \num{1.5} & \num{1.5} & \num{1.5} & \num{1.5} & \num{1.5} & \num{1.5} & \num{0.0} \\
{III.B Gasto neto pagado} & \num{26.2} & \num{27.0} & \num{25.5} & \num{24.9} & \num{24.4} & \num{24.3} & \num{24.2} & \num{24.1} & \num{-1.4} \\
{Gasto programable pagado} & \num{18.8} & \num{19.7} & \num{17.8} & \num{18.0} & \num{18.0} & \num{17.9} & \num{17.9} & \num{17.8} & \num{0.0} \\
{Diferimiento de pagos} & \num{-0.1} & \num{-0.1} & \num{-0.2} & \num{-0.1} & \num{-0.1} & \num{-0.1} & \num{-0.1} & \num{-0.1} & \num{0.1} \\
{Gasto programable devengado} & \num{18.9} & \num{19.8} & \num{18.0} & \num{18.1} & \num{18.1} & \num{18.0} & \num{18.0} & \num{17.9} & \num{-0.1} \\
{Gasto de operacion} & \num{15.7} & \num{16.1} & \num{15.3} & \num{15.4} & \num{15.7} & \num{15.6} & \num{15.6} & \num{15.6} & \num{0.3} \\
{Servicios personales} & \num{5.1} & \num{5.1} & \num{4.7} & \num{4.7} & \num{4.7} & \num{4.7} & \num{4.7} & \num{4.7} & \num{0.0} \\
{Otros de operacion} & \num{3.5} & \num{3.9} & \num{3.0} & \num{2.6} & \num{2.6} & \num{2.5} & \num{2.5} & \num{2.5} & \num{-0.5} \\
{Subsidios} & \num{2.7} & \num{2.7} & \num{3.0} & \num{3.5} & \num{3.7} & \num{3.7} & \num{3.6} & \num{3.5} & \num{0.5} \\
{Pensiones y jubilaciones} & \num{4.4} & \num{4.4} & \num{4.5} & \num{4.6} & \num{4.6} & \num{4.7} & \num{4.8} & \num{4.8} & \num{0.3} \\
{Gasto de capital} & \num{3.2} & \num{3.7} & \num{2.8} & \num{2.7} & \num{2.4} & \num{2.4} & \num{2.4} & \num{2.4} & \num{-0.4} \\
{Inversión física} & \num{2.7} & \num{3.0} & \num{2.3} & \num{2.4} & \num{2.0} & \num{2.1} & \num{2.0} & \num{2.0} & \num{-0.3} \\
{Inversión financiera} & \num{0.5} & \num{0.7} & \num{0.4} & \num{0.3} & \num{0.3} & \num{0.3} & \num{0.3} & \num{0.3} & \num{-0.1} \\
{Gasto no programable} & \num{7.5} & \num{7.4} & \num{7.7} & \num{6.9} & \num{6.5} & \num{6.5} & \num{6.3} & \num{6.3} & \num{-1.4} \\
{Costo financiero} & \num{3.7} & \num{3.6} & \num{3.8} & \num{3.2} & \num{2.8} & \num{2.8} & \num{2.7} & \num{2.7} & \num{-1.1} \\
{Participaciones} & \num{3.7} & \num{3.7} & \num{3.7} & \num{3.5} & \num{3.5} & \num{3.5} & \num{3.5} & \num{3.5} & \num{-0.2} \\
{Adefas} & \num{0.1} & \num{0.1} & \num{0.1} & \num{0.1} & \num{0.1} & \num{0.1} & \num{0.1} & \num{0.1} & \num{0.0} \\
{IV. Balance primario} & \num{-1.2} & \num{-1.4} & \num{0.6} & \num{0.5} & \num{0.5} & \num{0.4} & \num{0.4} & \num{0.4} & \num{-0.2} \\
{SHRFSP} & \num{48.8} & \num{51.4} & \num{51.4} & \num{51.4} & \num{51.4} & \num{51.4} & \num{51.4} & \num{51.4} & \num{0.0} \\
{SHRFSP interno} & \num{37.4} & \num{38.5} & \num{39.8} & \num{40.5} & \num{40.8} & \num{41.1} & \num{41.4} & \num{41.7} & \num{1.9} \\
{SHRFSP externo} & \num{11.4} & \num{12.9} & \num{11.6} & \num{10.9} & \num{10.6} & \num{10.3} & \num{10.0} & \num{9.7} & \num{-1.9} \\
{Deuda neta del Sector Publico} & \num{48.7} & \num{51.0} & \num{50.9} & \num{50.9} & \num{50.9} & \num{50.9} & \num{50.9} & \num{50.9} & \num{0.0} \\
{Deuda neta interna} & \num{37.1} & \num{37.9} & \num{39.2} & \num{40.0} & \num{40.4} & \num{40.8} & \num{41.1} & \num{41.5} & \num{2.3} \\
{Deuda neta externa} & \num{11.5} & \num{13.1} & \num{11.7} & \num{10.9} & \num{10.5} & \num{10.1} & \num{9.8} & \num{9.4} & \num{-2.2} \\
{Saldo historico de la deuda bruta del SPNF} & \num{52.9} & \num{54.9} & \num{54.9} & \num{54.9} & \num{54.9} & \num{54.9} & \num{54.9} & \num{54.9} & \num{0.0} \\
{Limite máximo de gasto corriente estructural} & \num{9.8} & \num{10.1} & \num{10.1} & \num{10.1} & \num{10.1} & \num{10.1} & \num{10.1} & \num{10.1} & \num{0.0} \\
\end{longtable}
\par\vspace{2pt}\footnotesize Fuente: elaboración propia del ITED con información de la SHCP, CGPE 2025, Anexo III.2, p. 85. El CGPE 2025 se publicó sin capa de texto; las cifras se leyeron de la página renderizada y se validaron por identidad contable.
\normalsize


El escenario oficial mantiene el saldo histórico en 51.4\% del PIB los seis años,
con requerimientos que bajan a 2.9 desde 2027 y un superávit primario que se
sostiene entre 0.4 y 0.6 puntos. Dentro del gasto, dos trayectorias merecen
registro: \textbf{pensiones y jubilaciones suben de 4.5 a 4.8\% del PIB} y el
costo financiero baja de 3.8 a 2.7. La segunda depende de que la tasa nominal
llegue a 5.5\% y se quede ahí.

Conviene anotar lo que el escenario \emph{no} incorpora, porque el propio cuadro
lo hace visible: la inversión física baja de 2.3 a 2.0\% del PIB y se queda ahí
seis años. \textbf{La estabilidad del acervo se sostiene, en el escenario oficial,
sobre una inversión pública permanentemente baja.}

\section{Comparaciones estructurales}

\emph{Costo financiero sobre impuestos}: 0.20 en 2022, 0.23 en 2023, 0.256 en 2024
y \textbf{0.262 en 2025}. Uno de cada cuatro pesos de impuestos se va en intereses.

\emph{Saldo histórico sobre impuestos, en años de recaudación}: 3.6, 3.4, 3.40 y
\textbf{3.51}. La deuda equivale a tres años y medio de toda la recaudación
tributaria, y el indicador empeora en 2025 pese a la consolidación, porque el
acervo crece en pesos mientras la recaudación crece menos.

\textbf{Las dos razones se mueven en la dirección contraria al mensaje del
paquete.} La consolidación es real en el flujo y todavía no se ve en el acervo.

\clearpage
